\section{Method: Evarb-0.2}
\label{sec:method}

% Figure 1 -- system architecture. Drawn from the implementation: memlab (harness, tools, scorer),
% evarb/{arbiter,policy,render,clients}.py, agentrun_v2.parser, scripts/run_p2.py.
% Explicit mm coordinates: nothing overlaps, nothing leaves the text block.
\begin{figure}[t]
\centering
\begin{tikzpicture}[font=\footnotesize, x=1mm, y=1mm]

% ---------------- row A (y=46): the agent's control flow
\node[blkplain, minimum width=18mm, minimum height=9mm]  (task)   at (12,46) {Task\\question};
\node[blk,      minimum width=26mm, minimum height=11mm] (agent)  at (44,46) {\textbf{LLM agent}\\[-1pt]\scriptsize tool loop};
\node[blk,      minimum width=26mm, minimum height=11mm] (submit) at (80,46) {\cond{submit\_answer}\\[-1pt]\scriptsize value / \cond{conflict} / abstain};
\node[blkacc,   minimum width=27mm, minimum height=11mm] (policy) at (116,46) {\textbf{Basis-gated}\\[-1pt]\textbf{policy} \scriptsize P1--P4};
\node[blkplain, minimum width=20mm, minimum height=9mm]  (final)  at (151,46) {Final\\answer};

% ---------------- row B (y=28): the gate and the single notice
\node[dec, minimum width=16mm] (gate)   at (128,28) {strength\\strong?};
\node[blkwarn, minimum width=20mm, minimum height=9mm] (notice) at (152,28) {one reversible\\notice};

% ---------------- row C (y=10): retrieval and arbitration
\node[store,    minimum width=18mm, minimum height=12mm] (store) at (12,10) {Memory\\store};
\node[blkplain, minimum width=25mm, minimum height=9mm]  (retr)  at (44,10) {\cond{memory\_search}\\[-1pt]\scriptsize hashed BoW, top 4};
\node[blkacc,   minimum width=26mm, minimum height=11mm] (arb)   at (84,10) {\textbf{Evarb-0.2 arbiter}\\[-1pt]\scriptsize status, basis, strength};
\node[blkplain, minimum width=25mm, minimum height=7mm]  (reg)   at (74,-4) {registry lookup};

% ---------------- control flow
\draw[flow] (task) -- (agent);
\draw[flow] (agent) -- (submit);
\draw[flow] (submit) -- (policy);
\draw[flow] (policy) -- (final);
\node[lbl, above=0.4mm of final] {\scriptsize to the frozen scorer};

% ---------------- retrieval and arbitration
\draw[flow] (38,40.5) -- node[lbl, left, pos=0.45, xshift=-0.5mm] {\scriptsize tool calls} (38,14.5);
\draw[flow] (38,30) -- (26,30) -- (26,-4) -- (reg.west);
\draw[flow] (store) -- (retr);
\draw[flowacc] (retr) -- (arb);
\draw[flowacc] (arb.north) -- (84,26) -- node[lbl, above, pos=0.62]
      {\scriptsize labelled records: status, basis, strength} (52,26) -- (52,40.5);
\draw[flowacc] (arb.east) -| node[lbl, pos=0.72, right, xshift=0.5mm] {\scriptsize resolution} (106,40.5);

% ---------------- the gate
\draw[flow] (124,40.5) -- (gate.north);
\draw[flowwarn] (gate) -- (notice);

ode[lbl, text=cEnf, anchor=south] at (139,31.5) {\scriptsize yes};
\draw[flowwarn] (notice.north) -- (152,62) -- node[lbl, above, pos=0.55]
      {\scriptsize the agent may resubmit once} (44,62) -- (agent.north);
\node[lbl, text=cNone, align=center, anchor=north] at (128,20)
  {\scriptsize no $\rightarrow$ advise:\\[-1pt]\scriptsize the submission stands};

\node[lbl, text=cEnf, align=left, anchor=north west] at (114,7)
  {\scriptsize \textbf{safe submission:} if the\\[-1pt]\scriptsize resubmission is unusable, the\\[-1pt]\scriptsize pre-notice answer stands};

\begin{scope}[on background layer]
\node[grp, fit=(store)(retr)(arb)(reg), inner sep=2.5mm,
      label={[lbl, text=cNone, anchor=north west, xshift=1mm]below left:\scriptsize retrieval and arbitration}] {};
\end{scope}
\end{tikzpicture}

\vspace{1mm}
\caption[AGENTRUN / Evarb system architecture]{%
\textbf{System architecture.} Everything the study varies lives in the two shaded components. The
\emph{arbiter} decides what the agent is told about a group of records --- a status, the evidence basis for that
status, and a strength. The \emph{policy} decides what happens when the agent's submission disagrees with the
arbitration: it may send at most one reversible notice, and it enforces only when the resolution's strength passes
the gate. The store, the retrieval function, the tool schemas, the prompts and the scorer are identical in all
twelve conditions. Advisory conditions keep the arbiter and disable the policy; ungated conditions widen the gate
to accept moderate strength as well as strong.\\
\emph{Accessible description:} a left-to-right control flow along the top (task, agent, submission, policy, final
answer); a lower band in which the memory store feeds search and search feeds the arbiter, which both returns
labelled records to the agent and sends its resolution to the policy; below the policy a strength gate either
issues one notice, which returns to the agent along the top of the figure, or lets the submission stand.\\
\emph{Source:} drawn from the implementation (\texttt{memlab}, \texttt{evarb/\{arbiter,policy,render\}.py},
\srcfile{scripts/run_p2.py}); no component is hypothetical.}
\label{fig:architecture}
\end{figure}


\subsection{System architecture}

\Cref{fig:architecture} shows the harness. The agent is a local causal LM with a Hermes-style tool template. It
has three tools: \cond{memory\_search}, a hashed bag-of-words retrieval returning the top four records; a registry
lookup that may fail by design; and \cond{submit\_answer}, the terminal action, which takes a value, the literal
\cond{conflict}, or an abstention.

Between retrieval and the agent sits the \textbf{arbiter}, and between the agent's submission and the recorded
answer sits the \textbf{policy}. Everything we vary lives in those two places: the arbiter decides \emph{what the
agent is told} about the records, and the policy decides \emph{what happens} when the agent's submission disagrees
with that. The store, the retrieval function, the tool schemas, the prompts and the scorer are identical in all
twelve conditions.

\subsection{Arbitration}

The arbiter receives the records of one (entity, attribute) group. Each record has an identifier, a value, a text,
a source and a possibly-missing \cond{recorded\_on} date. It returns a resolution kind (\cond{unchallenged},
\cond{resolved}, \cond{contradiction}, \cond{unresolved}), an \emph{evidence basis}, a \emph{strength}
(\cond{strong}, \cond{moderate}, \cond{none}), the current value if any, and a per-record status.

\paragraph{Clause-level cues.} Each record text is split into clauses at sentence ends, semicolons and dashes.
Within a clause the arbiter looks for cues of seven kinds --- correction, supersession, dismissal, retention,
unconfirmed, contradiction, and a hypothetical value --- and qualifies each by its scope: \emph{negated} if a
negator occurs among the four tokens before the cue; \emph{quoted} if the cue is in quotation marks or the clause
contains a reporting verb; \emph{hypothetical} if the clause contains a conditional or modal marker. A clause
\emph{refers to a rival record} if it mentions the rival's value (token-bounded, numbers compared numerically) or
uses an anaphor for an earlier value when exactly one rival value exists. Validity dates are read after ``as of'',
``effective'', ``valid from'', ``since'' and similar phrases; a date equal to the record's own value is never
treated as a validity date.

\paragraph{Evidence bases.} Every resolution carries one of ten evidence bases, and strength is a function of the
basis and its conditions (\Cref{fig:bases}). The scope qualifiers are what separate this from a lexical rule
list: a correction cue inside a negated, quoted or hypothetical clause does not license the ordering that the same
cue would license when asserted.

% Figure 2 -- the arbitration workflow, in the order implemented by evarb/arbiter.py (arbitrate()).
% All content is kept inside x in [8,148] so the picture sits squarely in the text block.
\begin{figure}[t]
\centering
\begin{tikzpicture}[font=\footnotesize, x=1mm, y=1mm,
  wfstep/.style={blk, minimum width=44mm, minimum height=7.6mm},
  wfout/.style={blkacc, minimum width=42mm, minimum height=7.6mm}]

\node[blkplain, minimum width=44mm, minimum height=7.6mm] (in) at (36,92)
  {records of one (entity, attribute) group};
\node[blkplain, minimum width=56mm, minimum height=9mm] (clause) at (36,80)
  {clause-level cue extraction\\[-1pt]\scriptsize 7 cue kinds $\times$ \{negated, quoted, hypothetical\},\\[-2pt]\scriptsize rival reference, validity dates};

\node[wfstep] (s1) at (36,67) {1\quad one distinct value?};
\node[wfstep, minimum height=9mm] (s2) at (36,56) {2\quad set aside records that do not\\[-1pt]establish their own value};
\node[wfstep] (s3) at (36,45) {3\quad negated pair / ``both stand''?};
\node[wfstep] (s4) at (36,35) {4\quad asserted referenced edges?};
\node[wfstep] (s5) at (36,25) {5\quad validity dates on all?};
\node[wfstep] (s6) at (36,15) {6\quad unique latest \cond{recorded\_on}?};
\node[wfstep] (s7) at (36,5)  {7\quad tie-break, in order};
\node[wfstep, fill=cWarnFill, draw=cEnf] (s8) at (36,-6) {8\quad missing dates, no text dates};

\node[wfout] (o1) at (108,67) {\cond{unchallenged}};
\node[wfout] (o2) at (108,56) {remainder current \scriptsize(moderate)};
\node[wfout] (o3) at (108,45) {\cond{contradiction}};
\node[wfout, minimum height=8.4mm] (o4) at (108,35) {winner, or \cond{unresolved}\\[-1pt]\scriptsize if inconsistent};
\node[wfout] (o5) at (108,25) {order by valid time \scriptsize(strong)};
\node[wfout, minimum height=8.4mm] (o6) at (108,15) {order by record time\\[-1pt]\scriptsize strong; moderate under doubt};
\node[wfout, minimum height=11mm] (o7) at (108,5)
  {\scriptsize unreferenced correction (mod.) $\to$ reported or\\[-1pt]\scriptsize hypothetical cue (\cond{unresolved}) $\to$ unique\\[-1pt]\scriptsize authoritative source (mod.) $\to$ \cond{contradiction}};
\node[wfout, minimum height=8.4mm] (o8) at (108,-6) {\cond{unresolved}\\[-1pt]\scriptsize strength \cond{none}};

\foreach \a/\b in {in/clause, clause/s1, s1/s2, s2/s3, s3/s4, s4/s5, s5/s6, s6/s7, s7/s8}
  \draw[flow] (\a) -- (\b);
\foreach \s/\o in {s1/o1, s2/o2, s3/o3, s4/o4, s5/o5, s6/o6, s7/o7, s8/o8}
  \draw[flowacc] (\s) -- (\o);

\node[blkwarn, minimum width=56mm, minimum height=11mm, align=center] at (101,-19)
  {\scriptsize \textbf{Every outcome carries a basis and a strength.}\\[-1pt]
   \scriptsize The policy enforces strong only (\cond{EBG}); strong $+$\\[-2pt]
   \scriptsize moderate (\cond{EBG-all}); never (advisory conditions).};
\end{tikzpicture}

\vspace{1mm}
\caption[Evarb-0.2 arbitration workflow]{%
\textbf{The arbitration workflow.} Steps are applied in this order, the first whose condition holds decides, and a
step that does not fire falls through to the next. Everything shown is implemented in \srcfile{evarb/arbiter.py};
nothing in the figure is proposed future work. The design choice this paper tests is not the ordering --- which is
close to standard bitemporal practice --- but the right-hand column: every outcome is accompanied by the
\emph{basis} on which it was reached and a \emph{strength}, and the downstream policy may enforce only what that
strength permits. Step 2 implements the benchmark's v2 semantics, under which a record that does not establish its
own value leaves the established value standing.\\
\emph{Accessible description:} a vertical chain of eight numbered tests, from the record group down to ``missing
dates''; each test has an arrow to the right into the outcome it produces; a box at the bottom states that
enforcement is gated on the strength attached to every outcome.\\
\emph{Source:} \srcfile{evarb/arbiter.py}; \srcfile{EVIDENCE_ARBITRATION_SPEC.md} \S3--\S4.}
\label{fig:workflow}
\end{figure}


\paragraph{Precedence.} Arbitration proceeds in a fixed order (\Cref{fig:workflow}). One distinct value is
\cond{unchallenged}. Otherwise the arbiter first \emph{sets aside records that do not establish their own value} ---
unconfirmed or proposed values, hypothetical changes, records stating that the rival remains valid, records saying
they do not replace the rival, and records whose own value appears only in reported, conditional or negated clauses
--- and if the remainder share one value, that value is current at moderate strength. Then: negated pairs or an
asserted ``both stand'' marker give a contradiction; asserted referenced textual edges give a winner if they are
consistent and \cond{unresolved} if not; validity dates present on every established record order them; otherwise
\cond{recorded\_on} orders them, with a guard that lowers strength when the winner's own text signals doubt. Ties
are broken, in order, by an unreferenced correction claimed by exactly one value, a reported or hypothetical cue
(\cond{unresolved}), a unique authoritative source, and otherwise a contradiction. Missing dates with no text dates
give \cond{unresolved}. The words ``update'' and ``revised'' never break a tie alone, because they are often
incidental.

The set-aside step encodes the benchmark's v2 semantics: \textbf{a record that does not establish its own value
leaves the established value standing.} These semantics were adopted after a blind annotation of the first
validation split showed the earlier gold for negated and reported claims to be inconsistent; that change was made
and frozen before the confirmatory data existed.

% Figure 3 -- the ten implemented evidence bases, grouped by the strength they can carry.
% Names, conditions and strengths are taken verbatim from EVIDENCE_ARBITRATION_SPEC.md section 3
% (and evarb/arbiter.py). No basis is renamed, added or omitted.
\begin{figure}[t]
\centering
\begin{tikzpicture}[font=\footnotesize, x=1mm, y=1mm,
  base/.style={draw=cBoxEdge, line width=0.4pt, rounded corners=1.2pt, align=left,
               inner sep=2.4pt, text width=46mm, font=\scriptsize, text=cInk, minimum height=9mm},
  hdr/.style={align=center, font=\scriptsize\bfseries, text=white, inner sep=2.2pt,
              rounded corners=1.2pt, minimum height=5.2mm, text width=48mm}]

% ---- column headers: the gate admits the left column only
\node[hdr, fill=cEBG]   (hs) at (26,86) {STRONG\quad\textnormal{\scriptsize the only column the gate admits}};
\node[hdr, fill=cAll, text=cInk] (hm) at (80,86) {MODERATE\quad\textnormal{\scriptsize enforced only by \cond{EBG-all}}};
\node[hdr, fill=cNone]  (hn) at (134,86) {NONE\quad\textnormal{\scriptsize never enforced}};

% ---- strong
\node[base, fill=cAccentFill] (b1) at (26,78)
  {\textbf{\cond{explicit\_correction}}\\asserted correction or dismissal referring to the rival, at different dates};
\node[base, fill=cAccentFill] (b2) at (26,65.5)
  {\textbf{\cond{explicit\_supersession}}\\asserted replacement or revision referring to the rival, at different dates};
\node[base, fill=cAccentFill] (b3) at (26,53)
  {\textbf{\cond{same\_day\_correction}}\\asserted correction, dismissal or supersession between same-date records, referring to the rival};
\node[base, fill=cAccentFill] (b4) at (26,40.5)
  {\textbf{\cond{valid\_time}}\\every established record states a validity date; the latest is unique};
\node[base, fill=cAccentFill] (b5) at (26,28)
  {\textbf{\cond{record\_time}}\\unique latest \cond{recorded\_on}};
\node[base, fill=cAccentFill] (b6) at (26,15.5)
  {\textbf{\cond{genuine\_contradiction}}\\same time, different values, no ordering evidence, no cue of any kind};
\node[base, fill=cAccentFill] (b7) at (26,3)
  {\textbf{\cond{negated\_correction}}\\negated correction with an asserted ``both stand / disagree'' marker};

% ---- moderate (demotions of the same bases, plus two of their own)
\node[base, fill=cWarnFill] (m1) at (80,78)
  {\textbf{\cond{explicit\_correction} / \cond{explicit\_supersession}}\\\emph{demoted} when the correcting record is the older one};
\node[base, fill=cWarnFill] (m2) at (80,65.5)
  {\textbf{\cond{same\_day\_correction}}\\\emph{demoted} when the correction is unreferenced and only breaks a same-time tie};
\node[base, fill=cWarnFill] (m3) at (80,53)
  {\textbf{\cond{record\_time}}\\\emph{demoted} when the winner's own text signals doubt (``proposed'', ``planned'', \ldots)};
\node[base, fill=cWarnFill] (m4) at (80,40.5)
  {\textbf{\cond{genuine\_contradiction}}\\\emph{demoted} when any cue is present};
\node[base, fill=cWarnFill] (m5) at (80,28)
  {\textbf{\cond{negated\_correction}}\\negated replacement without a marker, retention of the rival, or an unconfirmed change: that record is set aside};
\node[base, fill=cWarnFill] (m6) at (80,15.5)
  {\textbf{\cond{source\_authority}}\\same-time tie in which exactly one value comes from an authoritative source};
\node[base, fill=cWarnFill] (m7) at (80,3)
  {\textbf{\cond{quoted\_or\_hypothetical}}\\the record's own change or value is hypothetical: that record is set aside};

% ---- none
\node[base, fill=white] (n1) at (134,78)
  {\textbf{\cond{quoted\_or\_hypothetical}}\\the only ordering cue is reported speech $\rightarrow$ \cond{unresolved}};
\node[base, fill=white] (n2) at (134,65.5)
  {\textbf{\cond{insufficient\_evidence}}\\missing dates without text dates; cyclic or inconsistent cues $\rightarrow$ \cond{unresolved}};

\node[base, draw=none, fill=none, text width=50mm] (note) at (134,48)
  {\textbf{Ten bases, three strengths.} Six bases can carry more than one strength depending on their
   conditions, which is why gating cannot be replaced by a fixed list of ``trusted'' cue types: the same
   cue licenses enforcement in one context and only advice in another.};

\node[base, draw=cEnf, fill=cWarnFill, text width=50mm] (note2) at (134,22)
  {\textbf{What the gate costs.} Moving the gate from \emph{strong} to \emph{strong\,+\,moderate} ---
   that is, enabling the middle column --- raised task success by $9.7$\pp{} and the harmful-answer rate from
   3.1\% to 6.9\% (\Cref{sec:ablations}).};

\draw[decorate, decoration={brace, amplitude=3.5pt}, draw=cEBG, line width=0.6pt]
  (51.8,83.5) -- (51.8,-2.5);
\end{tikzpicture}

\vspace{1mm}
\caption[The ten evidence bases of Evarb-0.2]{%
\textbf{The ten evidence bases, grouped by the strength they can carry.} Names and conditions are verbatim from
the frozen specification; no basis is renamed, added or omitted. The left column is what gated enforcement acts on.
Six of the ten bases appear in more than one column because their strength depends on their conditions --- a
correction whose correcting record is the older one is demoted, an ordering by record date is demoted when the
winning record's own text signals doubt, and a contradiction is demoted when any cue is present. This
context-dependence is the reason the mechanism reports a basis and a strength rather than a verdict.\\
\emph{Accessible description:} a three-column layout headed STRONG, MODERATE and NONE, holding seven, seven and two
entries respectively; a brace marks the left column as the only one the gate admits.\\
\emph{Source:} \srcfile{EVIDENCE_ARBITRATION_SPEC.md} \S3; \srcfile{evarb/arbiter.py}.}
\label{fig:bases}
\end{figure}


\paragraph{What the agent sees.} Labels are attached to each retrieved record in the \cond{memory\_search} result:
a status, the evidence basis in words, a certainty (\cond{established}, \cond{likely}, \cond{not established},
mapped from strength) and, where relevant, \cond{superseded\_by}, \cond{conflicts\_with}, \cond{valid\_from} and a
note. A superseded record is shown as \cond{historical} when the question asks about an earlier time, because
whether a record is historical depends on the question rather than on the records. An order the evidence cannot
establish is shown as \cond{unresolved}, with the instruction to read the record texts --- never as
``conflicting'', because in the pilot small models followed wrong cautionary labels.

\subsection{The basis-gated policy}

The policy is evaluated when the agent submits, and applies only to \emph{anchored} groups: groups the agent
retrieved whose entity is named in the question. It never sees the gold answer, and model identity is not an input.
The first matching rule wins:

\begin{description}[leftmargin=2.2em,itemsep=1pt,topsep=2pt,font=\normalfont\sffamily]
\item[P3] no \cond{memory\_search} and no successful registry result $\rightarrow$ request evidence (search memory;
if nothing relevant, the previous submission stands);
\item[P1] current question, and the answer equals a superseded value and not the current one $\rightarrow$ enforce
if enforceable, else advise;
\item[P2] current question, status \cond{conflict}, and every anchored multi-value group is resolved $\rightarrow$
enforce if all are enforceable, else advise;
\item[P4] the answer equals a value of a contradictory group and there is no registry result $\rightarrow$ enforce
toward \cond{conflict} if enforceable, else advise;
\item[--] anything else, including every unresolved group $\rightarrow$ no action.
\end{description}

\paragraph{Gating} is the single switch this paper is about. A resolution is \emph{enforceable} in mode
\cond{gated} if its strength is strong; in mode \cond{all} if it is strong \emph{or} moderate; in mode \cond{none},
never. \cond{EBG} uses \cond{gated}; \cond{EBG-all} uses \cond{all}; the advisory conditions use \cond{none}.

\paragraph{Notices.} A run receives at most one notice. It names the records involved and states the evidence basis
truthfully --- a same-date correction is described as a correction, not as an ordering by date. The agent may
resubmit; the policy never changes an answer itself.

\paragraph{Safe submission.} If the submission after a notice is unparsable or missing, or the step budget runs
out, the final answer reverts to the pre-notice submission. This bounds the cost of a notice: it can turn one
parsable answer into another, but it cannot turn a parsable answer into a format or step-limit failure.

\subsection{Conditions}
\label{sec:conditions}

\begin{table}[t]
\centering
\caption{The twelve conditions. All twelve run in every cell. \cond{EBG} is the proposed condition; \cond{RT} is
the pre-registered primary comparator. ``Enforced'' means the policy may block or correct a submission that
disagrees with the arbitration; ``gated'' restricts that to resolutions of strong strength.}
\label{tab:conditions}
\footnotesize
\setlength{\tabcolsep}{3pt}
\begin{tabular}{@{}llllccl@{}}
\toprule
\textbf{Condition} & \textbf{Labels shown} & \textbf{Enforcement} & \textbf{Search} & \textbf{Safe} & \textbf{Role} \\
 & & & \textbf{first} & \textbf{submit} & \\
\midrule
\cond{B}   & none (memory only)      & ---                     & ---  & ---  & no-label baseline \\
\cond{G}   & metadata (record date)  & ---                     & ---  & ---  & advisory metadata \\
\cond{R}   & metadata                & every label             & yes  & no   & prior enforcement \\
\cond{T}   & temporal arbiter        & ---                     & ---  & ---  & advisory temporal labels \\
\cond{RT}  & temporal arbiter        & every label             & yes  & no   & \textbf{primary comparator} \\
\midrule
\cond{EB}  & evidence basis          & ---                     & ---  & ---  & labels only, advisory \\
\textbf{\cond{EBG}} & \textbf{evidence basis} & \textbf{gated (strong only)} & \textbf{yes} & \textbf{yes} & \textbf{proposed} \\
\cond{EBG-all}      & evidence basis  & strong \emph{and} moderate & yes & yes & ablation: the gate \\
\cond{EBG-unsafe}   & evidence basis  & gated                   & yes  & no   & ablation: safe submission \\
\cond{EBG-noR3}     & evidence basis  & gated                   & no   & yes  & ablation: search-before-submit \\
\midrule
\cond{LA}  & LLM arbiter             & ---                     & ---  & ---  & LLM arbiter, advisory \\
\cond{LG}  & LLM arbiter             & gated                   & yes  & yes  & LLM arbiter, EBG policy \\
\bottomrule
\end{tabular}
\end{table}


All twelve conditions run in every cell. They are constructed so that each comparison of interest differs in one
design element where possible: \cond{EBG} vs \cond{EB} isolates enforcement, \cond{EBG} vs \cond{EBG-all} isolates
gating, \cond{EBG} vs \cond{EBG-unsafe} isolates safe submission, \cond{EBG} vs \cond{EBG-noR3} isolates
search-before-submit, and \cond{EBG} vs \cond{LG} isolates \emph{who} arbitrates while holding the policy fixed.

The pre-registered primary contrast, \cond{EBG} $-$ \cond{RT}, is deliberately not of that kind: it compares the
proposed condition with the strongest prior enforcing condition and therefore changes the arbiter, the gating, the
notice wording and safe submission together. It cannot attribute an effect to any one of them. The ablations in
\Cref{sec:ablations} do the attribution, and they are exploratory.

\paragraph{Baselines.} \cond{G} and \cond{R} use metadata-only labels that order records by \cond{recorded\_on}
alone and label same-date disagreements as conflicting. \cond{T} and \cond{RT} use the earlier temporal arbiter:
it orders by text validity dates where every competing record has one, else by a unique latest \cond{recorded\_on},
breaking ties by a correction cue naming the rival's value, and otherwise reports no ordering evidence. It has no
negation, quotation or modality scope, and its labels carry no strength --- which is exactly what gating needs and
cannot get from it. \cond{R} and \cond{RT} enforce every label they produce and have no safe submission.
\cond{LA} and \cond{LG} use the agent's own model as arbiter, filling the same output schema from a prompt that
states the same evidence rules in plain language; unparsable or inconsistent output becomes \cond{unresolved} and
is counted as a format failure.
